\section*{Chapter \ref{ch:g7:comparative-advantage} --- Gains from Trade: Comparative Advantage}

\begin{solution}{exo:g7:comparative-advantage:1}
Leo is faster at cakes (1 hour against 2); Pia is faster at scarves (3 hours against 4).
\end{solution}

\begin{solution}{exo:g7:comparative-advantage:2}
The 4 hours of a shirt could bake $4 \div 1 = 4$ loaves: the \omterm{def:g7:scarcity-and-opportunity-cost:opportunity-cost}{opportunity cost} of a shirt
is 4 loaves.
\end{solution}

\begin{solution}{exo:g7:comparative-advantage:3}
The Netherlands $82.5\,\%$, the United States $11.0\,\%$. The United States has the larger
population (340 million against 18 million).
\end{solution}

\begin{solution}{exo:g7:comparative-advantage:4}
The North in cars (50 hours against 80); the South in rice (10 hours against 20).
\end{solution}

\begin{solution}{exo:g7:comparative-advantage:5}
(a) North: $50 \div 20 = 2.5$ tonnes of rice per car; South: $80 \div 10 = 8$ tonnes.
(b) The North has the \omterm{def:g7:comparative-advantage:comparative-advantage}{comparative advantage} in cars ($2.5 < 8$). A tonne of rice costs
the North $20 \div 50 = 0.4$ car and the South $10 \div 80 = 0.125$ car, so the South has
the \omterm{def:g7:comparative-advantage:comparative-advantage}{comparative advantage} in rice.
\end{solution}

\begin{solution}{exo:g7:comparative-advantage:6}
(a) Ana: $10 \div 1 = 10$ loaves and $10 \div 5 = 2$ chairs; Bo: $10 \div 2 = 5$ loaves
and $10 \div 5 = 2$ chairs. Together 15 loaves and 4 chairs. (b) A chair costs Ana
$5 \div 1 = 5$ loaves and Bo $5 \div 2 = 2.5$ loaves: Bo has the \omterm{def:g7:comparative-advantage:comparative-advantage}{comparative advantage} in
chairs. (c) Ana bakes $20$ loaves; Bo makes $20 \div 5 = 4$ chairs. Together 20 loaves and
4 chairs: 5 more loaves, the same chairs.
\end{solution}

\begin{solution}{exo:g7:comparative-advantage:7}
(a) $3 \times 2.2 = 6.6$ loaves. (b) Omar: 3 shirts and $6.6$ loaves against 3 and 6
without trade, a gain of $0.6$ loaf. Nadia: 4 shirts and $21 - 6.6 = 14.4$ loaves against
4 and 12, a gain of $2.4$ loaves. (c) Nadia gains more than at $2.5$ (where each gained
$1.5$); the total gain is still 3 loaves.
\end{solution}

\begin{solution}{exo:g7:comparative-advantage:8}
(a) No: with the 4 hours of a shirt, Omar can bake 2 loaves himself, more than the $1.5$
offered. (b) No trade takes place, and there are no \omterm{def:g7:comparative-advantage:gains-from-trade}{gains from trade}.
\end{solution}

\begin{solution}{exo:g7:comparative-advantage:9}
A country that trades freely, at low transport costs, and yet \omterm{def:g6:a-connected-world:export}{exports} the goods whose
\omterm{def:g7:scarcity-and-opportunity-cost:opportunity-cost}{opportunity cost} is higher at home than abroad, and \omterm{def:g6:a-connected-world:import}{imports} those whose \omterm{def:g7:scarcity-and-opportunity-cost:opportunity-cost}{opportunity cost}
is lower at home: for Japan after 1859, exporting goods that had been dear in Japan
before the opening.
\end{solution}

\begin{solution}{exo:g7:comparative-advantage:10}
(a) Brenland: $2 \div 4 = 0.5$ cloth per unit of grain; Kestria: $12 \div 6 = 2$ cloth.
(b) The grain cost Brenland 2 hours; weaving the cloth would have taken it 4 hours.
(c) The cloth cost Kestria 6 hours; growing the grain would have taken it 12 hours.
(d) Both: Brenland saves $4 - 2 = 2$ hours, Kestria $12 - 6 = 6$ hours.
\end{solution}

\begin{solution}{exo:g7:comparative-advantage:11}
(a) $1.5 \div 0.5 = 3$ loaves: unchanged. (b) Yes: the two \omterm{def:g7:scarcity-and-opportunity-cost:opportunity-cost}{opportunity costs} are still
3 and 2 loaves, so every ratio between 2 and 3 gives both a gain. (c) A shirt would cost
Nadia $2 \div 1 = 2$ loaves, the same as Omar: the range of ratios shrinks to the single
value 2, and there is no gain from trade.
\end{solution}

\begin{solution}{exo:g7:comparative-advantage:12}
(a) No: the model has two producers, each of whom gains; it says nothing about the
groups inside a country. (b) It leaves out that weavers and farmers are different people,
that changing jobs takes time and skills, and that some people lose their income while
others gain. (c) For example, the jobs and pay of Brenland's weavers and farmers, and the
prices of cloth and grain, before and after trade with Kestria grew.
\end{solution}

\begin{solution}{pb:g7:comparative-advantage:1}
\textbf{1.} Brenland: grain 2 hours, cloth 4 hours; Kestria: grain 12 hours, cloth 6
hours.
\textbf{2.} Brenland in both goods (2 against 12, 4 against 6).
\textbf{3.} Brenland: $60 \div 2 = 30$ grain and $60 \div 4 = 15$ cloth; Kestria: $60
\div 12 = 5$ grain and $60 \div 6 = 10$ cloth (millions of units).
\textbf{4.} 35 grain and 25 cloth.
\textbf{5.} Brenland $4 \div 2 = 2$ grain; Kestria $6 \div 12 = 0.5$ grain.
\textbf{6.} Brenland $2 \div 4 = 0.5$ cloth; Kestria $12 \div 6 = 2$ cloth.
\textbf{7.} Kestria in cloth ($0.5 < 2$), Brenland in grain ($0.5 < 2$).
\textbf{8.} Kestria: $120 \div 6 = 20$ cloth. Brenland: 5 cloth use $5 \times 4 = 20$
hours; the other 100 hours grow $100 \div 2 = 50$ grain. Total: 50 grain and 25 cloth,
15 grain more than in question 4, with the same cloth.
\textbf{9.} Kestria: $20 - 10 = 10$ cloth and 10 grain.
\textbf{10.} Brenland: $50 - 10 = 40$ grain and $5 + 10 = 15$ cloth.
\textbf{11.} Kestria: the same 10 cloth, and 10 grain instead of 5: a gain of 5 grain.
Brenland: the same 15 cloth, and 40 grain instead of 30: a gain of 10 grain.
\textbf{12.} $5 + 10 = 15$ grain, the extra production of question 8.
\textbf{13.} Kestria receives $10 \times 1.5 = 15$ grain: a gain of $15 - 5 = 10$ grain.
Brenland keeps $50 - 15 = 35$ grain: a gain of $35 - 30 = 5$ grain. Kestria gains more
now; the total gain is still 15.
\end{solution}
